\section{The duality-confinement membrane theorem}
\label{sec:master_theorem}

This section proves the confinement theorem.  Its input is a sequence
of domination records whose traces vanish; its output is the
almost-everywhere vanishing of the odd readout and, when the readout is
separating, confinement of the mass of \(\mu\) to the fixed locus.  In the
framework's companion papers a positive budget of this kind, bounding
a carrier-side quantity at every stage, is called a \emph{membrane}
\citep{TsiokosNeedleKiller2026}; the theorem's name records that it is
the step turning such budgets into confinement.  In the vocabulary of
\Cref{sec:framework}, it is the mathematical consequence of the
structural law \(\GamSDTC\): it does not produce the domination
records, it uses them.

\begin{theorem}[Duality-confinement membrane theorem]
\label{thm:duality_confinement:master-theorem}
Let \(\IOL=(X,J,\mu,\psi,Y,\Jiso)\) be an involutive object ledger with
\(\int_X\|\psim\|^2\,d\mu<\infty\).  Suppose there are positive
trace-class operators \(B_n\) on \(Y\) such that
\[
  \AX\preceq B_n\quad(n=1,2,\ldots)
  \qquad\text{and}\qquad
  \trace B_n\to0\quad(n\to\infty).
\]
Then \(\AX=0\) and \(\psim(x)=0\) for \(\mu\)-almost every \(x\).  If,
in addition, \(\psim\) is qualitatively separating, then \(Jx=x\) for
\(\mu\)-almost every \(x\), so the mass of \(\mu\) is confined to
\(\Fix(J)\); in the finite weighted case, every object of positive
weight is fixed by \(J\).

Conversely, if \(\psim(x)=0\) for \(\mu\)-almost every \(x\)---in
particular, if \(Jx=x\) for \(\mu\)-almost every \(x\)---then
\(\AX=0\), and the constant sequence \(B_n=0\) satisfies the
hypotheses.
\end{theorem}

\begin{proof}
By \Cref{lem:duality_confinement:trace-identity}, \(\AX\) is positive
and trace-class.  For each \(n\), (T4) applied to
\(0\preceq\AX\preceq B_n\) gives
\[
  0\leq\trace\AX\leq\trace B_n .
\]
The middle quantity does not depend on \(n\) and the right-hand side
tends to zero, so \(\trace\AX=0\).  By (T3), \(\AX=0\).  The remaining
direct assertions are
\Cref{thm:duality_confinement:separation-confinement}.  For the
converse, \(\psim\) vanishes on \(\Fix(J)\), so almost-everywhere
fixedness implies \(\psim=0\) almost everywhere; then
\(\trace\AX=\int_X\|\psim\|^2\,d\mu=0\), so \(\AX=0\) by (T3), and
\(\AX\preceq0\).
\end{proof}

The converse shows that, for a given ledger, trace-vanishing domination
records exist if and only if \(\psim\) vanishes almost everywhere, and
hence, for a separating readout, if and only if the mass is confined to
the fixed locus.  The hypothesis is therefore equivalent to the
conclusion, and the theorem does not make confinement easier to prove.
Its role is to fix the form in which confinement is to be established
when \(\AX\) is known only through a bridge
\(\AX\preceq\Kminus+E\) to operators computed from other data
(\Cref{sec:domination}): then trace bounds on \(\Kminus\) and \(E\),
rather than on \(\AX\) itself, are what must be proved.

Two features of the proof are worth isolating.  First, it uses \(\AX\)
only through four properties: \(\AX\) is positive; the trace is
monotone for the Loewner order; the trace of a positive operator is
nonnegative; and a positive operator of trace zero is zero.  Second,
the geometric conclusion enters only at the end, through the trace
identity and separation.  The first observation gives an abstract form
of the theorem.

\begin{remark}[Abstract form]
\label{rem:duality_confinement:abstract-cone}
Let \(C\) be a set with a distinguished element \(0\), a relation
\(\preceq\), a subset of \emph{positive} elements, and a map
\(\tau:C\to\mathbb R\) such that \(\tau(c)\geq0\) for every positive
\(c\), and \(a\preceq b\) implies \(\tau(a)\leq\tau(b)\).
Let \(a\in C\) be positive with \(\tau(a)=0\) only if \(a=0\).  If
\(a\preceq b_n\) for all \(n\) and \(\tau(b_n)\to0\), then \(a=0\): indeed \(0\leq\tau(a)\leq\tau(b_n)\to0\).  \Cref{thm:duality_confinement:master-theorem} is
the case of positive trace-class operators with the Loewner order and
the operator trace.
\end{remark}

\paragraph{Where the theorem is used.}
A companion paper \citep{TsiokosRHviaSDTC2026} considers the
involution \(s\mapsto1-\bar s\) on the zeros of the Riemann zeta
function.  It states, as an explicit hypothesis, a recognition source
supplying domination records for an odd readout of the zeros, and
applies the fixed-locus conclusion of the theorem to that ledger.  The
present paper uses nothing from that work.
